% ============================================================================
\section{System Architecture}
\label{sec:arch}
% ============================================================================

\subsection{Three terminals, one authority}
The brief asks for three control terminals: the main entrance, Lecture Room~1 and Lecture
Room~2 (\req{2.3}). Each one is a separate program with its own inputs and outputs:

\begin{description}[leftmargin=1.2em]
  \item[Entrance terminal (the \emph{hub}).] Reads the RFID readers, the keypad and the
        break-glass button, decides who may pass, drives the two door servos, the LCD, the buzzer
        and the LED, and keeps the occupancy ledger. It is the \textbf{single authority}: it is the
        only code that makes access decisions or changes who is inside.
  \item[Room terminals (\emph{room nodes} 1 and 2).] Each reads its own PIR and light sensor,
        decides its own lamp level and reports its state and a heartbeat to the hub. A room keeps
        controlling its lights even if the hub stops.
\end{description}

The terminals share no variables. They exchange \textbf{messages only}, over the node
\textbf{bus} (\cref{sec:sw-bus}). Two supporting tasks serve all terminals: the
\textbf{storage gatekeeper}, the only code that touches the SD card, and the
\textbf{serial CLI}.

\begin{figure}[H]
\centering
\begin{tikzpicture}
  % column x: inputs 0, hub 5.2, outputs 10.4 ; rows y top-down
  \node[hw]   (pc)    at (0,   4.2) {Serial terminal};
  \node[task] (cli)   at (5.2, 4.2) {\textbf{CLI task} (serial menu)};
  \node[hw]   (cards) at (0,   2.0) {Card buttons\\+ reader switch};
  \node[hw]   (keys)  at (0,   0.8) {4$\times$4 keypad};
  \node[hw]   (glass) at (0,  -0.3) {Break-glass};
  \node[task,minimum width=44mm,minimum height=30mm] (hub) at (5.2, 0.85)
    {\textbf{HUB task}\\(entrance, single authority)\\[2pt]
     \footnotesize access\_policy (pure)\\ \footnotesize auth\_fsm (pure)\\
     \footnotesize occupancy ledger};
  \node[hw]   (door)  at (10.4, 2.0) {2 door servos};
  \node[task] (ui)    at (10.4, 0.8) {\textbf{ui task}\\\footnotesize LCD, buzzer, LED};
  \node[task] (sto)   at (10.4,-1.9) {\textbf{STORAGE task}\\\footnotesize owns the SD card\\\footnotesize hash chain};
  \node[hw]   (sd)    at (10.4,-3.6) {microSD};
  \node[task] (r1)    at (2.9, -2.9) {\textbf{ROOM 1 task}\\\footnotesize room\_logic (pure)};
  \node[task] (r2)    at (7.0, -2.9) {\textbf{ROOM 2 task}\\\footnotesize room\_logic (pure)};
  \node[hw]   (r1hw)  at (2.9, -4.3) {PIR, LDR, lamp};
  \node[hw]   (r2hw)  at (7.0, -4.3) {PIR, LDR, lamp};

  \draw[arr,<->] (pc) -- (cli);
  \draw[arr,dashed] (cli) -- node[lbl,right]{reads services, sends commands} (hub);
  \draw[arr] (cards) -- node[lbl,above]{scan} (hub.west |- cards);
  \draw[arr] (keys)  -- node[lbl,above]{key} (hub.west |- keys);
  \draw[arr] (glass) -- (hub.west |- glass);
  \draw[arr] (hub.east |- door) -- (door);
  \draw[arr] (hub.east |- ui) -- (ui);
  \draw[arr,<->] (r1) -- node[lbl,left]{bus} (hub.south -| r1);
  \draw[arr,<->] (r2) -- node[lbl,right]{bus} (hub.south -| r2);
  \draw[arr,<->] (r1) -- (r1hw);
  \draw[arr,<->] (r2) -- (r2hw);
  \draw[arr,dashed] (hub.south east) -- node[lbl,sloped,above]{cocis\_log()} (sto.north west);
  \draw[arr] (sto) -- (sd);
\end{tikzpicture}
\caption{Runtime architecture. Blue boxes are FreeRTOS tasks, grey boxes are hardware. Solid arrows
between tasks are bus messages; every task logs through \code{cocis_log()}, which queues the line
to the storage task.}
\label{fig:arch}
\end{figure}

\subsection{Why one ESP32 board simulates the three terminals}
\label{sec:one-board}
The brief says each terminal \emph{may} use a separate ESP32 (\req{2.3}); it does not require it.
We run all three terminals on \textbf{one ESP32-S3} in Wokwi, built so that splitting them onto three
boards later changes only the transport:
\begin{enumerate}
  \item \textbf{Isolation is real, not cosmetic.} Each terminal is its own FreeRTOS task with its own
        pins and state. Room nodes never call hub functions and the hub never reads room variables;
        everything passes through \code{bus_send()} / \code{bus_recv()}.
  \item \textbf{The bus behaves like a network.} Messages carry a source id and a per-sender
        sequence number; the receiver drops duplicates and stale frames, the hub detects silent
        nodes by missing heartbeats, and a room can be ``unplugged'' with \code{node 2 offline} to
        demonstrate failure handling (\cref{sec:sync}).
  \item \textbf{Simulation practicality.} A single Wokwi project shows the whole building at once
        in one simulator, which makes the demonstration and the shared link reliable.
\end{enumerate}
On three physical boards, the queue transport in \file{components/bus/bus.c} would be replaced by
ESP-NOW or Wi-Fi frames carrying the same \code{bus_msg_t} structure; the terminal code above it
stays the same. This replacement transport is \emph{designed but not implemented}
(\cref{sec:limitations}).

\subsection{Layered components}
The firmware is split into ESP-IDF components with one concern each. Dependencies only point
downward, so lower layers never know about higher ones.

\begin{figure}[H]
\centering
\begin{tikzpicture}[layer/.style={box,minimum height=11mm,text width=#1,align=center}]
  % 3 columns of 34 mm, 3 mm gaps: centres at 0, 37, 74 mm; full width 108 mm
  \node[layer=104mm] at (37mm, 0)       {\textbf{main}: boot order};
  \node[layer=67mm]  at (18.5mm,-16mm)  {\textbf{nodes}: hub, rooms, door, rfid};
  \node[layer=30mm]  at (74mm, -16mm)   {\textbf{cli}: serial menu};
  \node[layer=30mm]  at (0,    -32mm)   {\textbf{ui}: LCD, sounds};
  \node[layer=30mm]  at (37mm, -32mm)   {\textbf{services}: users, config, occupancy, auth, clock};
  \node[layer=30mm]  at (74mm, -32mm)   {\textbf{bus}: node messages};
  \node[layer=30mm]  at (0,    -48mm)   {\textbf{drivers}: I2C, LCD, RTC, keypad, PWM, ADC};
  \node[layer=30mm]  at (37mm, -48mm)   {\textbf{storage}: SD gatekeeper, log};
  \node[layer=30mm,pure] at (74mm,-48mm){\textbf{core\_logic}: pure rules (unit tested)};
  \node[layer=104mm] at (37mm, -64mm)   {\textbf{common}: types, pin map, task table, SHA-256};
\end{tikzpicture}
\caption{Component layers (\code{firmware/components/}). The dashed box has no ESP-IDF calls, so it
compiles and is tested on a PC.}
\label{fig:layers}
\end{figure}

\begin{center}
\small
\begin{tabularx}{\linewidth}{@{}lL@{}}
\toprule
Component & Responsibility \\
\midrule
\code{common} & Shared enums (declared once with X-macros, which also generate their string
  tables), the pin map \file{board_pins.h}, the task table \file{task_config.h}, SHA-256. \\
\code{core_logic} & \textbf{Pure decision logic}: access rules, room lighting, admin two-factor state
  machine. No hardware, no RTOS, so it is unit-tested on a PC. \\
\code{drivers} & Thin hardware drivers: I2C bus, LCD2004, DS1307 RTC, keypad, buttons, PWM (servo,
  lamp, buzzer), light sensor. \\
\code{bus} & Messages between terminals with sequence numbers and duplicate rejection. \\
\code{storage} & The SD card gatekeeper task and the hash-chained log. \\
\code{services} & Shared state, each module behind its own mutex and persisted in NVS: users,
  config (schedule, lockdown), occupancy ledger, auth session, clock. \\
\code{ui} & The task that owns the LCD, buzzer and green LED. \\
\code{nodes} & The three terminals (hub, two rooms), the door servos and the RFID front end. \\
\code{cli} & The serial management menu (\code{esp_console}). \\
\bottomrule
\end{tabularx}
\end{center}
